% !PRINTED 478
\weberpage{478}
%<<CONT
minant $-4m$ is to be separated off, and if $m=1$, then the
equation remaining after the separation of the extraneous
factors is a square.

As examples of these cases we take:

1. \quad $m=3,\ n=3,$ \quad separation of a linear factor.\\
2. \quad $m=2,\ n=5,$ \quad no extraneous factor.\\
3. \quad $m=2,\ n=3,$ \quad separation of a quadratic factor.\\
4. \quad $m=1,\ n=3,$ \quad square.\\
5. \quad $m=1,\ n=5,$ \quad square after separation of a quadratic factor.\\
6. \quad $m=1,\ n=7,$ \quad square.

In case 1. one has to set, in the formula of § 73 relating to the
degree of transformation $n=3$:
\[
u^3=2,\quad v^3=f(\sqrt{-27})^3=2x,
\]
hence
\[
A=x^2+\frac{1}{x^2},\quad B=4x-\frac{2}{x},
\]
and consequently:
\[
x^4-4x^3+2x+1=(x-1)(x^3-3x^2-3x-1)=0,
\]
so that for $m=27$ one obtains:
\weq{1}{f(\sqrt{-27})^3=2x,\quad x^3-3x^2-3x-1=0.}

In case 2. we set, in system II. of § 73 ($n=5$):
\[
u_1=f_1(\sqrt{-2})=\sqrt[4]{2},\quad v_1=f_1(\sqrt{-50})=\sqrt[4]{2}\,x,
\]
\[
A_1=\frac{1}{x^3}-x^3,\quad B_1=2\Bigl(x^2+\frac{1}{x^2}\Bigr),
\]
so that (without an extraneous divisor) one obtains the equation of degree 6:
\weq{2}{\frac{1}{x^3}-x^3+2\Bigl(x^2+\frac{1}{x^2}\Bigr)=0,\quad f_1(\sqrt{-50})=\sqrt[4]{2}\,x}
which, for
\[
y=\frac{1}{x}-x
\]
reduces to the third degree:
\weq{3}{y^3+2y^2+3y+4=0.}
